\section{本章小结}
\label{sec:bilevel-summary}

双层优化的外层目标要读取一次优化的结果。本章首先分清三种对象：内解唯一时的$\Phi$，多解时靠规则确定的$\Phi_{\mathrm{o}}$和$\Phi_{\mathrm{p}}$，以及运行有限步算法得到的$\Phi_T$。乐观模型在紧性和连续性下有解，悲观模型没有同样的保证，算法的隐含择解不自动等于其中任何一种。内层无解的$\bm{x}$要排除，不能使用算法在那里给出的数字。

对解映射，Hessian正定给出局部的可微分支，全局强凸才让分支代表唯一的内解。驻点Jacobian可逆只保证有驻点分支，不保证是极小点；内解唯一也不保证可微。非凸内层的全局选解可能跳变，外层目标跟着不连续，算法沿分支算出的导数穿过切换点后就不再描述全局选解。

求导有两条路线。隐式微分在解点处解一个伴随方程，得到$\nabla\Phi=\nabla_{\bm{x}}F-\bm{B}^{\top}\bm{v}$，第二项是内解的响应，内层病态时线性残差被$1/\mu$放大；约束内层要解KKT系统，活动集变化、严格互补失效或KKT矩阵奇异时这个公式不能用。迭代展开对算法输出求导，得到的是$\Phi_T$的精确导数，反向累积省去大矩阵但要保存轨迹，截断反传则是导数的近似。

两条路线的关系由定理~\ref{thm:bilevel-unroll-limit}回答：内层一致收缩且二阶导数Lipschitz时，展开导数以$T\rho^T$的速率趋于隐式导数；只有点列收敛不够，固定的短展开也不是真实目标的梯度。

单层重述把内层最优性变成约束。KKT重述在凸内层和约束资格下全局等价，但产生的互补约束使MFCQ处处失效；值函数重述不需要凸性，代价是$V$可能不光滑、约束退化，松弛后要重新核对等价性。悲观模型不能这样改写。

表~\ref{tab:bilevel-summary}按误差来源归纳要核对的内容。外层收敛的曲线只回答了最后一行，其余各行需要各自的证据；局部驻定也不等于全局最优。

\begin{BookInlineTable}
 \centering
 \begin{tabular}{lll}
  \toprule
  误差来源 & 度量 & 控制手段\\
  \midrule
  内层求解 & $\delta$、$r$、$d$ & 增加内层步数\\
  内层选解 & 不同初值的内层值 & 重解，勿仅靠热启动\\
  伴随或截断 & $\lVert\bm{r}_{\bm{v}}\rVert_2$、$\ell$ & 增加线性迭代或反传长度\\
  外层 & 驻定残差、可行性 & 增加轮数，内层精度递增\\
  \bottomrule
 \end{tabular}
 \caption{双层求解中需要分别核对的四类误差。}
 \label{tab:bilevel-summary}
\end{BookInlineTable}
